\section{おわりに}
\label{sec:conclusion}

本研究は，工場内AGV環境において，
残存走行経路をAR風に可視化するeHMIに\textbf{危険度連動の視覚エンコーディング}を
導入する提案手法を設計した．
危険度は，視線方向の動的基準線とAGV残存経路のTTC，
および近接距離から統合的に算出される．

VR実験（被験者内3条件×10ケース）により，
移動効率とニアミスを定量的に比較し，
均一表示（Baseline）および非表示（No-AR）に対する
提案手法の優位性を検証する計画である．
シード固定によるAGV挙動統制，自動計測，構成概念の明示的操作化により，
実験計画は方法論的に妥当である．

今後，データ収集と統計分析を行い，
スマートファクトリーにおける人--AGV共存の安全HRI設計に
知見を還元する．
